% ==============================================================================
% Section 3: Trend Analysis: 2015-2024
% ==============================================================================

\section{Trend Analysis: 2015--2024}

A multi-year longitudinal analysis is crucial for discerning structural transformations from cyclical fluctuations in direct investment data. The decade from 2015 to 2024 witnessed momentous economic episodes: the abandonment of the CHF minimum exchange rate peg by the SNB in January 2015, the global COVID-19 pandemic shock, the global implementation of the OECD BEPS initiative, and the subsequent worldwide inflation and monetary tightening cycle.

\subsection{Inward FDI Capital Stocks Trend}

Between 2015 and 2024, Switzerland's inward FDI capital stock traversed a marked expansion followed by a protracted, steady contraction. Table~\ref{tab:inward_trend} chronicles the annual trajectory.

\begin{table}[H]
\centering
\small
\caption{Inward Foreign Direct Investment Capital Stocks in Switzerland (2015--2024)}
\label{tab:inward_trend}
\renewcommand{\arraystretch}{1.15}
\begin{tabularx}{0.85\textwidth}{c r r >{\raggedright\arraybackslash}X}
\toprule
\rowcolor{tableheadbg}
\textbf{Year} & \textbf{Total Stock ($\CHF$\,bn)} & \textbf{Annual Change (\%)} & \textbf{Trajectory Phase} \\
\midrule
2015 & 937.1 & --- & Base Period \\
\rowcolor{tablerowbg}
2016 & 1,264.0 & +34.9\% & Post-Peg Capital Influx \\
2017 & 1,347.8 & +6.6\% & All-Time Historical Peak \\
\rowcolor{tablerowbg}
2018 & 1,344.8 & -0.2\% & Plateau Phase \\
2019 & 1,345.0 & +0.0\% & Plateau Phase \\
\rowcolor{tablerowbg}
2020 & 1,239.5 & -7.8\% & Pandemic Disruption \& Tax Reform \\
2021 & 1,094.6 & -11.7\% & Supply Chain Realignments \\
\rowcolor{tablerowbg}
2022 & 1,033.7 & -5.6\% & Global Rate Hikes \& Valuation Shifts \\
2023 & 936.3 & -9.4\% & Cross-Border Deleveraging \\
\rowcolor{tablerowbg}
2024 & 925.7 & -1.1\% & Stabilization at Structural Baseline \\
\bottomrule
\end{tabularx}
\vspace{0.3em}
\footnotesize
\textit{Source: Swiss National Bank (SNB). Figures in billions of Swiss Francs ($\CHF\,$bn).}
\normalsize
\end{table}

\noindent\textbf{Key Empirical Observations on Inward Trends:}
\begin{itemize}[leftmargin=1.5em, itemsep=2.5pt]
    \item \textbf{The 2016--2017 Surge:} Inward FDI expanded rapidly from $\CHF\,937.1\bn$ in 2015 to an all-time peak of $\CHF\,1,347.8\bn$ in 2017 ($+43.8\pct$ increase), propelled by significant cross-border acquisitions and balance sheet expansions of European holding vehicles.
    \item \textbf{Cumulative Post-Peak Contraction:} From 2018 onwards, inward FDI contracted persistently. Between the 2017 peak and year-end 2024, cumulative inward stock dropped by \textbf{$\CHF\,422.1\bn$}, representing a \textbf{$31.3\pct$ reduction}.
    \item \textbf{Exchange Rate Valuation Deflation:} Because many foreign affiliates in Switzerland maintain parent-level accounting in Euros or US Dollars, the secular appreciation of the Swiss Franc mechanically compressed the recorded CHF book value of inward investment assets upon translation.
\end{itemize}

\subsection{Outward FDI Capital Stocks Trend}

Swiss outward direct investment exhibited significantly greater resilience and stability over the same ten-year period, reflecting the diversified international portfolio of Swiss multinational corporations (Table~\ref{tab:outward_trend}).

\begin{table}[H]
\centering
\small
\caption{Swiss Outward Foreign Direct Investment Capital Stocks Abroad (2015--2024)}
\label{tab:outward_trend}
\renewcommand{\arraystretch}{1.15}
\begin{tabularx}{0.85\textwidth}{c r r >{\raggedright\arraybackslash}X}
\toprule
\rowcolor{tableheadbg}
\textbf{Year} & \textbf{Total Stock ($\CHF$\,bn)} & \textbf{Annual Change (\%)} & \textbf{Trajectory Phase} \\
\midrule
2015 & 1,130.7 & --- & Base Period \\
\rowcolor{tablerowbg}
2016 & 1,362.2 & +20.5\% & Broad Global Expansion \\
2017 & 1,398.1 & +2.6\% & Steady Accumulation \\
\rowcolor{tablerowbg}
2018 & 1,511.1 & +8.1\% & Pre-Pandemic Investment High \\
2019 & 1,467.4 & -2.9\% & Moderate Asset Realignment \\
\rowcolor{tablerowbg}
2020 & 1,513.0 & +3.1\% & All-Time Historical Peak \\
2021 & 1,457.1 & -3.7\% & Post-Pandemic Rebalancing \\
\rowcolor{tablerowbg}
2022 & 1,312.0 & -10.0\% & Global Equity Market Corrections \\
2023 & 1,308.0 & -0.3\% & Consolidation Floor \\
\rowcolor{tablerowbg}
2024 & 1,340.1 & +2.5\% & Renewed MNE Outward Expansion \\
\bottomrule
\end{tabularx}
\vspace{0.3em}
\footnotesize
\textit{Source: Swiss National Bank (SNB). Figures in billions of Swiss Francs ($\CHF\,$bn).}
\normalsize
\end{table}

\noindent\textbf{Key Empirical Observations on Outward Trends:}
\begin{itemize}[leftmargin=1.5em, itemsep=2.5pt]
    \item \textbf{Peak at $\CHF\,1,513.0\bn$ in 2020:} Unlike inward FDI which peaked in 2017, outward direct investment reached its zenith in 2020, demonstrating the counter-cyclical strength and global liquidity of Swiss corporate giants during the initial COVID-19 pandemic shock.
    \item \textbf{Asset Valuation Correction in 2022:} The $10.0\pct$ contraction in 2022 was predominantly an equity market and bond price mark-to-market correction across North American and European subsidiaries rather than active physical disinvestment.
    \item \textbf{Rebound in 2024:} Outward capital stocks expanded by $+2.5\pct$ to $\CHF\,1,340.1\bn$ in 2024, confirming that Swiss MNEs have resumed strategic direct investments in pharmaceuticals, clean technology, and digital logistics.
\end{itemize}

% ------------------------------------------------------------------------------
% Pgfplots Figure 3.1: Inward vs Outward FDI Capital Stocks 2015-2024
% Monochromatic Greyscale Design with Distinct Markers and Legends
% ------------------------------------------------------------------------------
\begin{figure}[H]
\centering
\begin{tikzpicture}
\begin{axis}[
    width=14cm,
    height=7.2cm,
    grid=both,
    grid style={line width=0.2pt, color=midgray!30},
    xlabel={\small\textbf{Year}},
    ylabel={\small\textbf{FDI Capital Stock ($\CHF$\,Billion)}},
    xmin=2014.5, xmax=2024.5,
    ymin=800, ymax=1650,
    xtick={2015, 2016, 2017, 2018, 2019, 2020, 2021, 2022, 2023, 2024},
    xticklabel style={font=\footnotesize},
    yticklabel style={font=\footnotesize},
    legend style={
        at={(0.03,0.95)},
        anchor=north west,
        font=\footnotesize,
        fill=white,
        draw=pureblack,
        line width=0.6pt
    },
    line width=1.3pt
]

    % Outward FDI Curve (Pure Black with Filled Squares)
    \addplot[
        color=pureblack,
        mark=square*,
        mark size=2.2pt,
        line width=1.4pt
    ] coordinates {
        (2015, 1130.7)
        (2016, 1362.2)
        (2017, 1398.1)
        (2018, 1511.1)
        (2019, 1467.4)
        (2020, 1513.0)
        (2021, 1457.1)
        (2022, 1312.0)
        (2023, 1308.0)
        (2024, 1340.1)
    };
    \addlegendentry{Swiss Outward FDI Stock (Abroad)}

    % Inward FDI Curve (Dark Gray with Dashed Line and Filled Circles)
    \addplot[
        color=darkgraytext,
        dashed,
        mark=*,
        mark size=2.2pt,
        line width=1.4pt
    ] coordinates {
        (2015, 937.1)
        (2016, 1264.0)
        (2017, 1347.8)
        (2018, 1344.8)
        (2019, 1345.0)
        (2020, 1239.5)
        (2021, 1094.6)
        (2022, 1033.7)
        (2023, 936.3)
        (2024, 925.7)
    };
    \addlegendentry{Inward FDI Stock (in Switzerland)}

    % Annotate Net Capital Exporter Gap in 2024
    \draw[{Stealth[length=2mm, width=1.5mm]}-{Stealth[length=2mm, width=1.5mm]}, line width=1pt, color=pureblack] 
        (axis cs:2024, 935) -- (axis cs:2024, 1330)
        node[midway, left=2pt, font=\scriptsize\bfseries, align=right] {Net Capital Export Surplus\\$\Delta = \CHF\,414.4\bn$};

    % Peak Annotation
    \node[anchor=south, font=\tiny\bfseries] at (axis cs:2017, 1355) {Inward Peak: 1,347.8};
    \node[anchor=south, font=\tiny\bfseries] at (axis cs:2020, 1520) {Outward Peak: 1,513.0};

\end{axis}
\end{tikzpicture}
\caption{Comparative Evolution of Swiss Inward versus Outward Direct Investment Capital Stocks (2015--2024). Source: SNB. The persistent positive vertical distance confirms Switzerland's structural role as a net capital exporter.}
\label{fig:trend_comparison}
\end{figure}

\subsection{FDI Integration as Percentage of GDP}

To gauge the relative intensity of direct investment within the broader economy, direct investment capital stocks are benchmarked against Switzerland's nominal Gross Domestic Product (GDP), which stood at approximately \textbf{$\CHF\,800\bn$} in 2024:

\begin{equation}
    \text{Inward FDI Intensity} = \frac{\text{Inward FDI Stock}_{2024}}{\text{Nominal GDP}_{2024}} = \frac{\CHF\,925.7\bn}{\CHF\,800.0\bn} \approx \mathbf{115.7\%}
\end{equation}

\begin{equation}
    \text{Outward FDI Intensity} = \frac{\text{Outward FDI Stock}_{2024}}{\text{Nominal GDP}_{2024}} = \frac{\CHF\,1,340.1\bn}{\CHF\,800.0\bn} \approx \mathbf{167.5\%}
\end{equation}

By comparison, the average OECD inward and outward direct investment intensity ratios fluctuate between $40\%$ and $55\%$ of GDP. Switzerland's combined direct investment ratio ($\approx 283\%$ of GDP) places it in an elite tier of hyper-globalized investment hubs alongside Singapore, Ireland, and Luxembourg, demonstrating unprecedented exposure to and leadership within the global division of labor.
